\documentclass[11pt]{article}
\usepackage[margin=1.1in]{geometry}
\usepackage{amsmath,amssymb,amsthm}
\usepackage[colorlinks=true,linkcolor=blue,citecolor=blue]{hyperref}

\newtheorem{theorem}{Theorem}[section]
\newtheorem{lemma}[theorem]{Lemma}
\newtheorem{proposition}[theorem]{Proposition}
\newtheorem{corollary}[theorem]{Corollary}
\newtheorem{definition}[theorem]{Definition}
\newtheorem{remark}[theorem]{Remark}
\newtheorem{problem}[theorem]{Problem}

\newcommand{\R}{\mathbb{R}}
\newcommand{\E}{\mathbb{E}}
\newcommand{\PP}{\mathbb{P}}
\newcommand{\emp}{\bar\mu_N}
\newcommand{\vol}{\operatorname{vol}}
\newcommand{\diam}{\operatorname{diam}}

\title{The critical scaling of the uniform doubling condition:\\
smooth bodies versus truncations}
\author{Yitzchak Shmalo\\
\small Einstein Institute of Mathematics, The Hebrew University of Jerusalem, Jerusalem, Israel\\
\small }
\date{June 2026}

\begin{document}
\maketitle

\begin{abstract}
The uniform doubling condition of [Shmalo, \emph{Ann.\ Math.\ Artif.\ Intell.}
(2024)] is a property of a training set $T\subset\R^n$ guaranteeing stability of
accuracy during the training of neural networks with piecewise linear
activations: for every truncated slab and every width on a geometric grid between
a scaling factor $\ell$ and a scale parameter $\beta$, widening the slab by the
factor $\kappa$ must multiply its empirical mass by $1+\sigma$ unless the mass
already exceeds $\delta$.  In the companion paper [\emph{The uniform doubling
condition holds with high probability for independent samples}] we proved that
$N$ i.i.d.\ samples with density bounded above and below on a convex body $K$
satisfy the condition with probability $1-1/N$, with $N$-independent
$(\kappa,\delta,\sigma,\beta)$ and scaling factor
$\ell_N=O\bigl((\log N/N)^{1/n}\bigr)$, and that some power-law decay of
$\ell_N$ is necessary: on the ball, the condition fails with probability bounded
below unless $\ell\ge c\,N^{-2/(n+1)}$.  This left a gap between the exponents
$1/n$ and $2/(n+1)$.  Here we close the gap and reveal a dichotomy.  For
\emph{pure slabs} (no truncations) over a body with an interior rolling ball, a
depth-integral estimate
$\vol\bigl(K\cap\{|\langle y-x,u\rangle|\le L\}\bigr)\ge c\,L^{(n+1)/2}$, valid
uniformly in $x\in K$ and in the direction $u$, improves the scaling factor to
$\ell_N=O\bigl((\log N/N)^{2/(n+1)}\bigr)$: the critical exponent for smooth
bodies is exactly $2/(n+1)$, up to logarithms.  For \emph{truncated} slabs the
improvement is impossible: already one fixed, uniformly thick truncation of the
cube creates, with probability bounded below, a diagonal-level gap beyond the
cornermost sample point which forces $\ell\ge c\,N^{-1/n}$.  Thus the truncation
number, not the smoothness of the support, dictates the critical scaling:
$2/(n+1)$ without truncations, $1/n$ with them.
\end{abstract}

\section{Introduction}\label{sec:intro}

Let $K\subset\R^n$ be a convex body, let $\mu$ be a probability measure with
density $\rho$ on $K$ satisfying
\begin{equation}\label{eq:density}
0<\rho_{\min}\le\rho\le\rho_{\max}<\infty\ \text{ on }K,
\qquad\widetilde R:=\rho_{\max}/\rho_{\min},
\end{equation}
and let $T=\{X_1,\dots,X_N\}$ be i.i.d.\ samples from $\mu$ with empirical
measure $\emp$.  For a unit vector $u$, a center $s\in\R^n$, a width $w>0$ and a
truncation set $h=\{(v_i,t_i)\}_{i\le k}$, the truncated slab is
\[
S_w\;=\;\Bigl\{x:\ |\langle x-s,u\rangle|\le\tfrac w2\Bigr\}\cap H,
\qquad
H=\{x:\ \langle v_i,x\rangle\le t_i,\ i\le k\},
\]
with truncation region $G:=H\cap K$.  The \emph{uniform doubling condition}
(UDC) of \cite{ShmaloUDC}, with constants $(\kappa,\delta,\sigma,\beta,\ell)$ and
truncation family $\mathcal L$, requires
\[
\emp(S_{\ell\kappa^{i+1}})\;\ge\;
\min\bigl\{\delta,\ (1+\sigma)\,\emp(S_{\ell\kappa^{i}})\bigr\}
\]
for every truncated slab generated by $\mathcal L$ and every
$i\in\mathbb N_0$ with $\ell\kappa^i<\beta$.  The condition is the hypothesis of
the stability-of-accuracy theorems of \cite{BJS,ShmaloUDC}.  In the companion
paper \cite{companion2} we proved:

\begin{theorem}[{\cite{companion2}}]\label{thm:old}
Under \eqref{eq:density}, for every $k$ and $r_0>0$ there are
$N$-independent constants $\kappa=1+150\widetilde R$, $\sigma=1$,
$\delta=\frac{3\rho_{\min}\omega_nr_0^n}{32\widetilde R}$, $\beta=\diam K$, such
that with probability at least $1-\frac1N$ the sample satisfies the UDC,
simultaneously for all truncation families of size $k$ whose nonempty regions
$G$ contain a ball of radius $r_0$, with scaling factor
\[
\ell_N\;=\;C_0\Bigl(\frac{\widetilde R}{\rho_{\min}}
\Bigl(\frac{\diam K}{r_0}\Bigr)^{n-1}\frac{\log N}{N}\Bigr)^{1/n}.
\]
Moreover, for the uniform measure on the unit ball and untruncated slabs, with
probability bounded below by a constant the condition fails for every
$\ell\le c\,N^{-2/(n+1)}$ {\rm(}for any fixed $\kappa>1$, $\sigma>0$, and
$\delta\ge2/N${\rm)}.
\end{theorem}

The exponents do not match: $1/n<\frac{2}{n+1}$ for $n\ge2$, i.e.\ the
sufficient scaling factor is larger than the necessary one.  This paper
determines the critical exponent, with an answer that depends not on the
regularity of $K$ but on the presence of truncations.

\begin{theorem}[Smooth bodies, pure slabs]\label{thm:smooth}
Let $K$ have the interior rolling ball property at radius $\rho_K>0$: for every
boundary point $p$ there is a ball $B\subseteq K$ of radius $\rho_K$ with
$p\in\partial B$ {\rm(}true, e.g., for any $K$ with $C^2$ boundary and strictly
positive principal curvatures; for the unit ball, $\rho_K=1${\rm)}.  Assume
\eqref{eq:density} and consider the UDC for \emph{untruncated} slabs ($k=0$,
$G=K$).  Then with probability at least $1-\frac1N$ the sample satisfies the
UDC with the constants $(\kappa,\sigma,\delta,\beta)$ of Theorem~\ref{thm:old}
{\rm(}with $r_0=\rho_K${\rm)} and scaling factor
\[
\ell_N\;=\;C_1\,\Bigl(\frac{\log N}{N}\Bigr)^{2/(n+1)},
\qquad
C_1=C_1\bigl(n,\widetilde R,\rho_{\min},\rho_K\bigr).
\]
Together with the lower bound of Theorem~\ref{thm:old}, the critical exponent
for bodies with a rolling ball and pure slabs is $\frac{2}{n+1}$, up to the
logarithmic factor.
\end{theorem}

\begin{theorem}[Truncations shift the exponent]\label{thm:wedge}
Let $K=[0,1]^n$ with $\mu=\mathrm{Leb}$, $n\ge1$, and fix any
$\kappa>1$, $\sigma>0$, $\beta>0$ and $\delta\ge\frac3N$.  Consider the single
truncation set $h^*=\{(-\mathbf 1/\sqrt n,\ -(n-\tfrac12)/\sqrt n)\}$, i.e.
\[
G^*\;=\;\Bigl\{x\in[0,1]^n:\ \textstyle\sum_jx_j\ge n-\tfrac12\Bigr\},
\]
a corner cap containing a ball of radius $r_0(n)=\frac{1}{4n}$.  There are
constants $c_2(n,\kappa),c_3(n,\kappa)>0$ such that for all large $N$, with
probability at least $c_3$, the doubling inequality fails for a truncated slab
generated by $h^*$ at \emph{every} grid scale $\le c_2N^{-1/n}$; hence the UDC
with truncation number $k\ge1$ forces
\[
\ell\;\ge\;c_2\,N^{-1/n}.
\]
Consequently the scaling factor of Theorem~\ref{thm:old} is sharp up to
logarithms in the truncated setting --- even for a single, fixed, uniformly
thick truncation --- and the critical exponent is $\frac1n$ with truncations
versus $\frac{2}{n+1}$ without.
\end{theorem}

The mechanism is geometric.  The empirical argument of \cite{companion2}
requires that any (truncated) slab of width $w\ge\ell_N$ containing a sample
point have, after widening by $\kappa$, population mass of order at least
$\log N/N$, so that the Vapnik--Chervonenkis relative deviation bounds take
over.  Without truncations, the widened slab around \emph{any} point of a
rolling-ball body captures at least the volume of a spherical-cap stack,
$\asymp L^{(n+1)/2}$ in the half-width $L$
(Lemma~\ref{lem:rolling}); equating with $\log N/N$ gives the exponent
$\frac2{n+1}$.  With truncations, the slab may live inside a corner region of
$G$ whose level sections degenerate polynomially, where the captured mass is
only $\asymp L^n$; the matching lower-bound construction places the slab at the
cornermost sample point, beyond which a level gap of relative size $\kappa^2$
occurs with probability bounded below, stalling the empirical mass at $1/N$
through an entire grid step.

In the regime relevant to deep networks --- where the truncations are the
activation cells, $k\ge1$ --- the exponent $\frac1n$ of \cite{companion2} is
therefore correct; for the pure-slab doubling conditions of
Berlyand--Jabin--Safsten type \cite{BJS} on smooth data supports, the faster
scaling $\frac2{n+1}$ is available.

\section{Ingredients from the companion paper}\label{sec:prelim}

We reuse three components of \cite{companion2}.

\begin{proposition}[{Population dichotomy, \cite{companion2}}]
\label{prop:population}
Let $\kappa\ge4$ and let $S_w$ be a truncated slab whose region $G$ contains a
ball of radius $r_0$.  Then for every $w>0$,
\[
\mu(S_{\kappa w})\;\ge\;\min\Bigl\{\delta_0,\
\Bigl(1+\tfrac{\kappa-1}{15\widetilde R}\Bigr)\mu(S_w)\Bigr\},
\qquad
\delta_0=\frac{3\rho_{\min}\omega_nr_0^n}{16\widetilde R}.
\]
\end{proposition}

\begin{lemma}[{VC relative deviations, \cite{companion2}}]
\label{lem:vc}
Let $d=d_{n,k}$ be the VC bound for intersections of $k+2$ half-spaces and
$u_N=\frac4N\bigl(d\ln\frac{2eN}d+\ln(8N)\bigr)$.  With probability at least
$1-\frac1N$, simultaneously for all truncated slabs $C$:
if $\mu(C)\ge16u_N$ then $\emp(C)\ge\frac34\mu(C)$; and if
$\emp(C)\ge4u_N$ then $\mu(C)\ge\frac12\emp(C)$.
\end{lemma}

\begin{proposition}[{Proof skeleton, \cite{companion2}}]\label{prop:skeleton}
Let $N$ be large enough that $32u_N\le2\delta_0$.  On the event of
Lemma~\ref{lem:vc}, suppose that a scale $w$ has the property that every
truncated slab $S$ of width $w$ (region $G\supseteq$ some $r_0$-ball) satisfies
\begin{equation}\label{eq:keymass}
\emp(S_w)>0
\;\Longrightarrow\;
\mu(S_{\kappa w})\;\ge\;32\,u_N .
\end{equation}
Then the doubling inequality
$\emp(S_{\kappa w})\ge\min\{\delta,2\emp(S_w)\}$, with
$\delta=\delta_0/2$ and $\kappa=1+150\widetilde R$, holds at scale $w$ for all
such slabs.
\end{proposition}

Indeed, inspection of the proof of the main theorem of \cite{companion2} shows
that the sample point $x\in S_w\cap G$ is used only to produce the mass bound
\eqref{eq:keymass} (there, through a cone with apex $x$ and base an
$r_0$-ball, giving mass $\gtrsim w^n$); every subsequent step --- the case
analysis combining Proposition~\ref{prop:population} with
Lemma~\ref{lem:vc} --- consumes \eqref{eq:keymass} alone.  The scaling factor
$\ell_N$ may therefore be taken as any $w_0$ such that \eqref{eq:keymass} holds
for all $w\in[w_0,\beta)$.  Improving the geometric lower bound improves the
exponent, and this is all we shall do.

\section{Pure slabs over rolling-ball bodies}\label{sec:rolling}

\begin{lemma}[Depth-integral slab-volume bound]\label{lem:rolling}
Let $K$ have the interior rolling ball property at radius $\rho_K$, and let
$n\ge2$.  Then for every $x\in K$, every unit vector $u$, and every
$0<L\le\rho_K$,
\[
\vol\Bigl(K\cap\bigl\{y:\ |\langle y-x,u\rangle|\le L\bigr\}\Bigr)
\;\ge\;
a_n\,\rho_K^{(n-1)/2}\,L^{(n+1)/2},
\qquad
a_n:=\frac{\omega_{n-1}}{2^{\,n}},
\]
where $\omega_{n-1}$ is the volume of the unit ball of $\R^{n-1}$.  (For $n=1$
the left-hand side is at least $L$, and the statement holds trivially with
$a_1=1$.)
\end{lemma}

\begin{proof}
Write $\rho:=\rho_K$, $\tau(y):=\langle y,u\rangle$,
$\tau_{\max}:=\max_K\tau$, $\tau_{\min}:=\min_K\tau$, and let
$A(\tau):=\vol_{n-1}\bigl(K\cap\{\tau(y)=\tau\}\bigr)$, so that the slab volume
equals $\int_{\tau(x)-L}^{\tau(x)+L}A(\tau)\,d\tau$.  Note that $K$ contains a
ball of radius $\rho$, so $\tau_{\max}-\tau_{\min}\ge2\rho\ge2L$.  By the
Brunn--Minkowski inequality, $A^{1/(n-1)}$ is concave on
$[\tau_{\min},\tau_{\max}]$; in particular $A$ is \emph{unimodal}, so on any
subinterval $A$ is bounded below by the smaller of its endpoint values.

\emph{Step 1: an anchor level near each extreme.}  Let $p\in\partial K$ attain
$\tau_{\max}$ and let $B=B(c,\rho)\subseteq K$ be a rolling ball with
$p\in\partial B$; set $g:=\tau(p)-\tau(c)\in[0,\rho]$ (nonnegative since
$\tau(c)\le\tau_{\max}$, and at most $|p-c|=\rho$).  The section of $B$ at the
level $\tau_{\max}-h$ is an $(n-1)$-ball of radius $r(h)$ with
$r(h)^2=\rho^2-(g-h)^2$.  For $0\le h\le g$,
\[
r(h)^2=(\rho-g+h)(\rho+g-h)\;\ge\;h\cdot\rho,
\]
since $\rho-g+h\ge h$ and $\rho+g-h\ge\rho$.  Hence
\begin{equation}\label{eq:section}
A(\tau_{\max}-h)\;\ge\;\omega_{n-1}(\rho h)^{(n-1)/2}
\qquad(0\le h\le g),
\end{equation}
and at $h=g$ the section of $B$ is its full equator:
$A(\tau_{\max}-g)\ge\omega_{n-1}\rho^{\,n-1}$.

Define the \emph{anchor depth} $a:=\min(g,\,L/4)$ and the value
$A^*:=\omega_{n-1}(\rho L/4)^{(n-1)/2}$.  Then
\begin{equation}\label{eq:anchor}
A(\tau_{\max}-a)\;\ge\;A^* :
\end{equation}
if $g\ge L/4$ then $a=L/4\le g$ and \eqref{eq:section} gives
$A(\tau_{\max}-a)\ge\omega_{n-1}(\rho L/4)^{(n-1)/2}=A^*$; if $g<L/4$ then
$a=g$ and the equator bound gives
$A(\tau_{\max}-a)\ge\omega_{n-1}\rho^{n-1}\ge A^*$, using $\rho\ge L/4$.
Symmetrically, with $g'$ the offset of a rolling ball at a minimizer of
$\tau$, the anchor depth $a':=\min(g',L/4)$ satisfies
$A(\tau_{\min}+a')\ge A^*$.  Note $a,a'\le L/4$.

\emph{Step 2: the integral.}  By \eqref{eq:anchor} and unimodality,
\[
A(\tau)\;\ge\;A^*
\qquad\text{for all }
\tau\in M:=[\tau_{\min}+a',\ \tau_{\max}-a].
\]
The complement of $M$ in $[\tau_{\min},\tau_{\max}]$ has total length
$a+a'\le L/2$.  The slab's level interval
$[\tau(x)-L,\tau(x)+L]$ intersects $[\tau_{\min},\tau_{\max}]$ in an interval
of length at least $L$: indeed $\tau(x)$ lies in the range, and the range has
length at least $2\rho\ge2L$, so at least one of the two sides of $\tau(x)$
contains a full segment of length $L$ inside the range.  Consequently the slab
meets $M$ in a set of length at least $L-\frac L2=\frac L2$, and
\[
\int_{\tau(x)-L}^{\tau(x)+L}A
\;\ge\;\frac L2\,A^*
\;=\;\frac{\omega_{n-1}}{2\cdot4^{(n-1)/2}}\,\rho^{(n-1)/2}L^{(n+1)/2}
\;=\;\frac{\omega_{n-1}}{2^{\,n}}\,\rho^{(n-1)/2}L^{(n+1)/2}.
\]
\end{proof}

\begin{proof}[Proof of Theorem \ref{thm:smooth}]
Take $r_0=\rho_K$ (the rolling ball is contained in $K$, so the inradius is at
least $\rho_K$, and $G=K$ trivially contains an $r_0$-ball).  Set
\[
\ell_N:=\frac{2}{\kappa-1}
\Bigl(\frac{32\,u_N}{\rho_{\min}\,a_n\,\rho_K^{(n-1)/2}}\Bigr)^{2/(n+1)},
\]
with $u_N$ from Lemma~\ref{lem:vc} for $k=0$, and let $N$ be large enough that
$32u_N\le\min\{2\delta_0,\ \rho_{\min}a_n\rho_K^{n}\}$ and
$\frac{(\kappa-1)\ell_N}{2}\le\rho_K$.  By
Proposition~\ref{prop:skeleton} it suffices to verify \eqref{eq:keymass} for
every $w\ge\ell_N$ and every slab $S_w$ ($k=0$).  Let $x\in S_w\cap T$ be a
sample point (so $x\in K$ a.s.), and put $L:=\frac{(\kappa-1)w}{2}$.  Since
$|\langle x-s,u\rangle|\le\frac w2$, the slab
$\{y:|\langle y-x,u\rangle|\le L\}$ is contained in $S'_{\kappa w}(u,s)$, so by
Lemma~\ref{lem:rolling} (applied with half-width $\min(L,\rho_K)$),
\[
\mu(S_{\kappa w})\;\ge\;\rho_{\min}\,a_n\,\rho_K^{(n-1)/2}\,
\min(L,\rho_K)^{(n+1)/2}.
\]
If $L\le\rho_K$, the right-hand side is at least
$\rho_{\min}a_n\rho_K^{(n-1)/2}\bigl(\tfrac{(\kappa-1)\ell_N}2\bigr)^{(n+1)/2}
=32u_N$ by the definition of $\ell_N$ and $w\ge\ell_N$.  If $L>\rho_K$, it is
at least $\rho_{\min}a_n\rho_K^{n}\ge32u_N$ by the choice of $N$.  In both
cases \eqref{eq:keymass} holds, and the theorem follows with
$C_1$ determined by $u_N=O\bigl(\tfrac{n\log N}{N}\bigr)$ for $k=0$.
\end{proof}

\section{Truncations: the corner-gap obstruction}\label{sec:wedge}

\begin{proof}[Proof of Theorem \ref{thm:wedge}]
Write $e:=\mathbf 1/\sqrt n$, $\tau(x):=\langle x,e\rangle$, and
$\tau_{\max}=\sqrt n$, attained at the corner $\mathbf 1$.  The region
$G^*=\{x\in K:\tau(x)\ge\tau_{\max}-\eta_0\}$ with
$\eta_0:=\frac{1}{2\sqrt n}$ is a simplex cap at the corner.  It contains a
ball of radius $r_0=\frac1{4n}$: the point
$z:=\bigl(1-\tfrac{\eta_0}{2\sqrt n}\bigr)\mathbf 1$ has
$\tau(z)=\sqrt n\,\bigl(1-\tfrac{\eta_0}{2\sqrt n}\bigr)
=\tau_{\max}-\tfrac{\eta_0}{2}$; each coordinate of $z$ equals
$1-\tfrac{\eta_0}{2\sqrt n}$, so $B(z,r)\subseteq[0,1]^n$ for
$r\le\tfrac{\eta_0}{2\sqrt n}=\tfrac1{4n}$; and on $B(z,r)$ one has
$\tau\ge\tau(z)-r\ge\tau_{\max}-\eta_0$ since $r\le\tfrac{\eta_0}2$.  Hence
$B\bigl(z,\tfrac1{4n}\bigr)\subseteq G^*$.

For $h\in[0,\eta_0]$ let
$\Lambda(h):=\mu\bigl(\{x\in K:\ \tau(x)\ge\tau_{\max}-h\}\bigr)$ be the
corner-cap mass.  The cap $\{\tau\ge\tau_{\max}-h\}$ in the cube is the corner
simplex $\{\sum_j(1-x_j)\le h\sqrt n\}$, of volume
$\frac{(h\sqrt n)^n}{n!}$ for $h\sqrt n\le1$; thus
\begin{equation}\label{eq:cap}
\Lambda(h)=\frac{n^{n/2}}{n!}\,h^{n}
\qquad\text{for }0\le h\le\tfrac1{\sqrt n}.
\end{equation}

Fix $\kappa,\sigma,\delta\ge\frac3N,\beta$ and set
$\Lambda_0:=\frac{n^{n/2}}{n!}$, $h_N:=\bigl(\frac{1}{\Lambda_0N}\bigr)^{1/n}$
and $h^*:=4\kappa^2h_N$, so $\Lambda(h_N)=\frac1N$ and
$\Lambda(h^*)=\frac{(4\kappa^2)^n}{N}$.  Consider the event
\[
E:=\Bigl\{\text{exactly one sample lies in the cap }
\{\tau\ge\tau_{\max}-h^*\},\ \text{and it lies in }
\{\tau\ge\tau_{\max}-h_N\}\Bigr\}.
\]
As in the binomial computation of the sharpness section of \cite{companion2},
\[
\PP(E)\;=\;N\,\Lambda(h_N)\bigl(1-\Lambda(h^*)\bigr)^{N-1}
\;\ge\;N\cdot\frac1N\Bigl(1-\frac{(4\kappa^2)^n}{N}\Bigr)^{N-1}
\;\ge\;c_3(n,\kappa)\;>\;0
\]
for all large $N$.  Work on $E$ and let $X^*$ be the unique sample point in
the $h^*$-cap, with depth $h_1:=\tau_{\max}-\tau(X^*)\le h_N$.

Consider the truncated slabs $S_w$ with direction $u:=e$, center $s:=X^*$, and
truncation set $h^*$ (region $G^*$); every such slab contains $X^*$, so
$\emp(S_w)\ge\frac1N>0$ for all $w>0$.  On the other hand, since $S_w\subseteq
G^*\cap\{\tau\ge\tau(X^*)-\frac w2\}\subseteq\{\tau\ge\tau_{\max}-h_1-\frac
w2\}$ and $h_1\le h_N=\frac{h^*}{4\kappa^2}\le\frac{h^*}4$, every slab of
width $w\le h^*$ satisfies
$S_w\subseteq\{\tau\ge\tau_{\max}-h^*\}$, the cap containing no sample other
than $X^*$ on the event $E$.  Hence for all $w\le h^*/\kappa$,
\[
\emp(S_w)\;=\;\emp(S_{\kappa w})\;=\;\frac1N
\;<\;\min\Bigl\{\delta,\ \frac{1+\sigma}{N}\Bigr\}
\;=\;\min\Bigl\{\delta,\ (1+\sigma)\,\emp(S_w)\Bigr\},
\]
because $\delta\ge\frac3N>\frac1N$ and $\sigma>0$.  Thus the doubling
inequality is violated at every scale $w\le h^*/\kappa$.

If $\ell\le h^*/\kappa^2$, the grid scale $w=\ell\kappa^{i}$ with $i=0$
already lies in $(0,h^*/\kappa]$ (and $\ell<\beta$ for large $N$), so the UDC
fails on $E$.  Hence on $E$ the UDC forces
$\ell>h^*/\kappa^2=4h_N=c_2(n,\kappa)N^{-1/n}$.
\end{proof}

\begin{remark}
The example uses the corner of the cube only for the explicit cap volume
\eqref{eq:cap}; any $G$ whose level sections in some direction degenerate like
$h^{n-1}$ --- equivalently, any truncation region with a conical point of its
closure in the relevant direction, no matter how thick the region is globally
--- produces the same obstruction.  Smoothness of $K$ is irrelevant: the
truncation hyperplanes create the corners.  This is why the truncated exponent
$\frac1n$ cannot be improved by regularity assumptions on the data support
alone.
\end{remark}

\section{Endpoint questions}\label{sec:final}

\begin{problem}
Determine the critical scaling at the logarithmic level: does the UDC for pure
slabs on the ball hold with $\ell_N\asymp N^{-2/(n+1)}$ (no logarithm), or is
a logarithmic factor necessary?  The upper bound consumes
$\log N$ through the VC bound of Lemma~\ref{lem:vc}; the lower-bound
construction of \cite{companion2} produces only a single violating slab, and a
union over $\asymp N^{c}$ essentially independent boundary caps should improve
it to $\ell\ge c\,(\log N/N)^{2/(n+1)}$, matching the upper bound.  We leave
the details to future work.
\end{problem}

\begin{problem}
Interpolate between the two regimes: for truncation regions whose boundary
consists of smooth pieces meeting along edges of dimension $j$
{\rm(}$0\le j\le n-1${\rm)}, determine the critical exponent of the scaling
factor as a function of the lowest edge dimension present in the truncation
family.  The two theorems of this paper are the extreme cases: rolling-ball
boundary (no edges) gives $\frac{2}{n+1}$, and a conical point ($j=0$) gives
$\frac1n$; we expect the level-section degeneracy at a $j$-dimensional edge to
produce an intermediate exponent, monotone in $j$.
\end{problem}

\subsection*{Acknowledgments}
The research presented in this paper was supported by the European Research
Council (ERC) under the European Union's Horizon Europe research and innovation
programme (grant agreement No.\ 101041711), by the Israel Science Foundation
(grant number 2258/19) and by the Israel Science Foundation (ISF Grant
4101/25).

\subsection*{Declaration on generative AI}
During the preparation of this work the author used generative AI tools to
accelerate drafting and revision.  All mathematical claims, proofs, and
bibliographic information were subsequently reviewed and edited by the author,
who takes full responsibility for the content of the manuscript.

\begin{thebibliography}{9}

\bibitem{BJS}
L.~Berlyand, P.-E.~Jabin and C.~A.~Safsten,
\emph{Stability for the training of deep neural networks and other classifiers},
Math.\ Models Methods Appl.\ Sci.\ \textbf{31} (2021), no.~11, 2345--2390.

\bibitem{ShmaloUDC}
Y.~Shmalo,
\emph{Stability of accuracy for the training of DNNs via the uniform doubling
condition},
Ann.\ Math.\ Artif.\ Intell.\ (2024), DOI 10.1007/s10472-023-09919-1.

\bibitem{companion2}
Y.~Shmalo,
\emph{The uniform doubling condition holds with high probability for
independent samples},
manuscript, 2026.

\bibitem{Gardner}
R.~J.~Gardner,
\emph{The Brunn--Minkowski inequality},
Bull.\ Amer.\ Math.\ Soc.\ (N.S.) \textbf{39} (2002), no.~3, 355--405.

\end{thebibliography}

\end{document}
